
\def \Subject {شبکه‌های بیزی}
\def \Session { یک}
% \setcounter{chapter}{\Session-1}
\renewcommand{\chaptername}{سوال}
\chapter{\Subject}
\section{بخش اول}
\subsection{ زیربخش اول}

\begin{figure}[H]
	\centering
	\begin{tikzpicture}[
		node distance=1.5cm,
		var/.style={circle, draw, minimum size=0.8cm, align=center, font=\bfseries},
		arrow/.style={-latex, thick}  % Simple, no length parameter
		]
		
		% Define nodes
		\node[var] (W) {W};
		\node[var, right=of W] (E) {E};
		\node[var, below=of $(W)!0.5!(E)$] (P) {P};  % Requires calc library
		\node[var, below=of P] (A) {A};
		\node[var, below left=of A] (M) {M};
		\node[var, below right=of A] (O) {O};
		\node[var, below=of $(M)!0.5!(O)$] (H) {H};  % Requires calc library
		
		% Draw arrows
		\draw[arrow] (W) -- (P);
		\draw[arrow] (E) -- (P);
		\draw[arrow] (P) -- (A);
		\draw[arrow] (A) -- (M);
		\draw[arrow] (A) -- (O);
\draw[arrow, bend left=25] (P) to (H);  % Curved arrow
		\draw[arrow] (M) -- (H);
		\draw[arrow] (O) -- (H);
		
	\end{tikzpicture}
	\caption{شبکه بیزی آلودگی هوا}
	\label{fig:bayesian_network}
\end{figure}
\subsection{زیربخش دوم}
با توجه به شبکه رسم شده در شکل 
\ref{fig:bayesian_network}
داریم:
\begin{latin}
	~\\
	$
	P(X_1, X_2, X_3, ...) = \prod\limits_{1}^{n} P(X_i | parent_i)
	\\
	\Rightarrow
	P(W, E, P, A, M, O, H) = P(W) P(E) P(P | W, E) P(A | P) P(M | A) P(O | A) P(H | M, O, P)
	$
\end{latin}
\subsection{زیربخش سوم}
سوال 
\lr{Markov Blanket}
مربوط به متغیر 
\lr{M}
را می‌خواهد. پس داریم:
\begin{latin}
	~\\
	$
	MB(M) = Parents(M) + Children(M) + Spouse(M) 
	\\
	Parents(M) = {A} 
	\quad 
	Children(M) = {H}
	\quad 
	Spouse(M) = {P, O}
	\\
	\Rightarrow 
	MB(M) = {A, H, P, D}
	$
\end{latin}
\subsection{زیربخش چهارم}
باید درستی و یا نادرستی عبارات را تعیین کرد.
\begin{enumerate}
	\item 
	$
	A \perp  W | P
	$

بله درست است. در شبکه بیزی، 
\lr{A}
و
\lr{W}
تنها از طریق 
\lr{P}
باهم رابطه دارند. پس اگر مقدار 
\lr{P}
مشخص باشد، این دو از هم مستقل خواهند بود.
\item 
$
M \perp O | A
$

درست است. چون 
\lr{A}
تنها والد مشترک آن‌ها است، با مشخص بودن 
\lr{A}
دیگر متغیری وجود ندارد که بتواند هردوی متغیرهای 
\lr{M}
و
\lr{O}
را به‌هم وابسته کند. و خود متغیرها نیز با همدیگر رابطه‌ای ندارند.

\item 
$
A \perp H | \{M, O\}
$

غلط است. زیرا 
\lr{A}
و
\lr{H}
والد مشترک 
\lr{P}
را نیز همچنان دارند و این مورد روی هردو آن‌ها تاثیر می‌گذارد که یعنی حتی با معلوم بودن
\lr{M}
و
\lr{O}
نیز، این دو متغیر از هم مستقل نیستند.
\item 
$
M \perp E | A
$
درست است. مسیر دیگری بین دو متغیر 
\lr{E}
و
\lr{M}
نیست که از 
\lr{A}
گذر نکند.

\item 
$
O \perp W | \{A, H\}
$
غلط است. زیرا با مشخص بودن 
\lr{H}
عملا رابطه هم‌بستگی بین والدهای آن بوجود می‌آید و 
\lr{O}
والد مستقیم و 
\lr{W}
والد غیر مستقیم آن است که مسیر رابطه بین این‌ها از 
\lr{A}
عبور نمی‌کند.
\end{enumerate}

\section{بخش دوم}
باید مقایسه کرد که در هر مورد، کدام بزرگ‌تر است. 
\begin{enumerate}
	\item 
	$P(w^1) \qquad p(w^1 | f^1)$
	
	مقدار 
	$p(w^1)$
	بیشتر است. زیرا 
	$f^1$
	باعث کم شدن احتمال 
	$e^1$
	می‌شود و 
	$e^1$
	باعث زیاد شدن 
	$w^1$
	می‌شود. پس با رخ دادن 
	$f^1$ 
	احتمال 
	$w^1$
	کاهش می‌یابد. 
	
		\item 
	$P(f^1 | e^1) \qquad p(f^1 | e^1, h^1)$
	
	برابرند. تنها مسیر ارتباط 
	$f$
	با باقی از طریق $e$ است و با مشخص بودن 
	$e$
	دیگر مقدار سایر متغیرها تاثیری در احتمال 
	$f$ 
	نخواهد داشت. 
	\item 
	$P(h^1 | f^0, w^1) \qquad P(h^1 | w^1)$
	
	مقدار 
	$P(h^1 | w^1)$
	بیشتر است. زیرا می‌توان دید که احتمال 
	$w$
	با بالا رفتن احتمال 
	$h$
	زیاد می‌شود. حال اگر مقدار 
	$w$
	را داشته باشیم، اگر 
	$f$
	منفی باشد، ممکن است $f^0$ باعث مثبت شدن 
	$w$ 
	شود و برای همین به دلیل پدیده توضیح‌دهی، احتمال 
	$h^1$
	وقتی 
	$f^0$
	باشد کاهش می‌یابد.
	
	\item 
	$p(e^1 | w^1) \qquad p(e^1 | w^1, c^1)$
	
	مقدار $ p(e^1 | w^1, c^1)$ بیشتر است. زیرا 
	$c^1$
	باعث کم شدن احتمال 
	$d^1$
	می‌شود و این به دلیل پدیده توضیح‌دهی باعث زیاد شدن احتمال
	$e^1$
	می‌شود. 
	\item 
	$p(w^1 | d^1, e^1) \qquad p(w^1 | d^1, e^1, h^1)$
	
	برابرند. زیرا همه مسیرهای رسیدن به 
	\lr{w}
	در هر دو حالت کاملا مشخص شده‌اند. 
	
	\item 
	$p(t^1 | w^0, e^1) \qquad p(t^1 | w^0, e^1, f^1)$
	
	برابرند. زیرا 
	$f$
	نمی‌تواند بدون تغییر 
	$e$
	اثری روی 
	$t$
	بگذارد. پس با مشخص بودن 
	$e$
	عملا $f$ بی‌تاثیر خواهد بود.
	
\end{enumerate}

\section{بخش سوم}
ما n متغیر داریم که هرکدام L مقدار مختلف می‌توانند داشته باشند. 
\subsection{زیربخش اول}
ابتدا به ازای n های کوچک حساب می‌کنیم:
\begin{latin}
	~\\
	$
	n = 1 \Rightarrow L - 1 \\
	n = 2 \Rightarrow L^2 - 1 \\
	n = 3 \Rightarrow L^3 - 1 
	$
\end{latin}
چون در هرمورد 
$L^n$
متغیر داریم. ولی چون جمع آن‌ها یک است، پس یکی از آن‌ها وابسته است و 
$L^n - 1$
متغیر مستقل داریم. 

\subsection{زیربخش دوم}
باید تعداد حالات مستقل خود گره را در تعداد حالات والدهای آن ضرب کرد. در بخش قبل حساب کردیم که 
n
متغیر، 
$L^n  - 1$
حالت مستقل دارند. البته در اینجا چون جمع این احتمالات ثابت نیست، منهای یک نداریم.
خود گره نیز 
$L - 1$
حالت دارد 
)
پس درنهایت یک گره حداکثر
$(L^k) ( L - 1)
$
پارامتر دارد. حال باید این عدد را در تعداد گره‌ها ضرب کرد پس درنهایت داریم:

\begin{latin}
	$
	\text{number of parameters} = \text{number of nodes} * \text{each node parent parametrs} * \text{each node number of states} = n * (L^k) * (L - 1) 
	$
\end{latin}
\subsection{زیربخش سوم}
ایده  شبیه بخش قبل است و فقط شیوه محاسبه فرق می‌کند. پس داریم:
\begin{latin}
	~\\
	$
	\text{number of parameters} = \sum\limits_{i = 0}^n param(X_i) + param(C)
	= \sum\limits_{i = 0}^n states(parent(X_i))states(X_i - 1) + (m - 1)
	= m - 1 + m * (L - 1) + m * L * (L - 1) + m * l * (L - 1) + ...
	= (m - 1) + (m * (L - 1))  + (n - 1)(m * L * (L - 1))
	= m - 1 + ml - m + ((n - 1) (mL^2 -mL ))
	= ml - 1 + nmL^2 - nmL - mL^2 + mL
	= 2-nml - 1 + nmL^2 - nmL - mL^2 
	$
\end{latin}
\section{بخش چهارم}
\subsection{زیربخش اول}
در هرمورد باید دید آیا یالی وجود دارد که با برعکس کردن آن هیچ 
\lr{v structure}
ای حذف نشود و هیچ 
\lr{v structure}
جدیدی بوجود نیاید. 

در گراف سمت چپ، با تغییر هر یک از یال‌های 
A-B
و 
C-B
باعث می‌شود 
\lr{V structure}
مربوط به 
A-B-C
خراب شود. از طرف دیگر  تغییر جهت یال‌های 
B-D
و
D-E
منجر به ایجاد 
\lr{V structure}
های جدید می‌شود. پس هیچ گراف دیگری با گراف داده شده 
\lr{I equivalent }
نیست. 

در گراف سمت راست، با تغییر جهت یال 
E-D
هیچ 
\lr{v strucutre}
ای حذف یا تولید نمی‌شود. پس این گراف با گراف سوال 
\lr{I equivalent }
است. شکل این گراف در تصویر
\ref{q1:4-equiv}
رسم شده است.
\begin{figure}[H]
	\centering
	\begin{tikzpicture}[
		node distance=1.5cm,
		var/.style={circle, draw, minimum size=0.8cm, align=center, font=\bfseries},
		arrow/.style={-latex, thick}  % Simple, no length parameter
		]
		
		% Define nodes
		\node[var] (A) {A};
		\node[var, right=of A] (B) {B};
		\node[var, right=of B] (C) {C};
		\node[var, right=of C] (D) {D};
		\node[var, right=of D] (E) {E};
		

		% Draw arrows
		\draw[arrow] (A) -- (B);
		\draw[arrow] (B) -- (C);
		\draw[arrow] (D) -- (E);
		\draw[arrow] (D) -- (C);

		
	\end{tikzpicture}
	\caption{گراف معادل گراف سمت راست سوال}
	\label{q1:4-equiv}
\end{figure}
\subsection{زیربخش دوم}
دو شرط آمده در سوال بیان می‌کنند که 
\lr{A}
و
\lr{B}
نسبت به هم مستقلند و همین‌طور هر دو والد مشترک \lr{C} نیستند. این دو شرط در کنارهم باعث می‌شود تا گراف همبند نشود. زیرا \lr{C} نمی‌تواند هم به 
\lr{A}
و
هم به 
\lr{B}
وصل باشد (فارغ از جهت) و 
A
و
B
نیز رابطه‌ای باهم ندارند.  پس \lr{A} و \lr{B} درون دو مولفه جدا از هم قرار می‌گیرند. ولی C در هرکدام از این مولفه‌ها باشد، به این معنی است که از متغیر دیگر مستقل است و این استقلال در صورت سوال نیامده است. پس نمی‌توان شبکه‌ای یافت که 
\lr{perfect map}
این توزیع باشد.
